% META: {"id": "TNSI-ECRIT-S5-EX1", "chapitre": "TNSI-STRUCTURES-DONNEES", "type_objet": "banque_ecrite", "sujet": "TNSI-ECRIT-S5", "rang": 1, "capacites": ["T-STRUCT-01A", "T-STRUCT-02B"], "capacites_codes": ["C1", "C5"], "duree_min": 70, "autorite": "MENE2516123N", "mode_creation": "ex_nihilo", "sources_inspiration": [], "notice": "ORIGINAL_NEXUS_TRAINING_MATERIAL_NOT_OFFICIAL_EXAM_CONTENT", "status": "generated"}
% BEGIN-VERIFY
% class Stock:
%     """Interface : ajouter, retirer, quantite, references, est_vide."""
%     def __init__(self):
%         self._quantites = {}
%     def ajouter(self, reference, nombre):
%         assert nombre > 0, "ajout non strictement positif"
%         self._quantites[reference] = self._quantites.get(reference, 0) + nombre
%     def retirer(self, reference, nombre):
%         assert nombre > 0, "retrait non strictement positif"
%         disponible = self._quantites.get(reference, 0)
%         assert nombre <= disponible, "stock insuffisant"
%         if nombre == disponible:
%             del self._quantites[reference]
%         else:
%             self._quantites[reference] = disponible - nombre
%     def quantite(self, reference):
%         return self._quantites.get(reference, 0)
%     def references(self):
%         return sorted(self._quantites)
%     def est_vide(self):
%         return self._quantites == {}
% # Q1 : l'interface tient ses promesses
% s = Stock()
% assert s.est_vide() and s.references() == [] and s.quantite("cle") == 0
% s.ajouter("cle", 4)
% s.ajouter("vis", 10)
% s.ajouter("cle", 3)
% assert s.quantite("cle") == 7
% assert s.references() == ["cle", "vis"]
% s.retirer("vis", 10)
% assert s.references() == ["cle"]           # la reference epuisee disparait
% assert s.quantite("vis") == 0
% # Q2 : les preconditions sont refusees a l'entree
% for appel in (lambda: s.ajouter("cle", 0), lambda: s.ajouter("cle", -1),
%               lambda: s.retirer("cle", 0), lambda: s.retirer("cle", 99),
%               lambda: s.retirer("absent", 1)):
%     refuse = False
%     try:
%         appel()
%     except AssertionError:
%         refuse = True
%     assert refuse
% assert s.quantite("cle") == 7              # aucun appel refuse n'a modifie l'etat
% # Q3 : une seconde implementation, meme interface
% class StockListe:
%     def __init__(self):
%         self._lignes = []                  # liste de [reference, nombre]
%     def _index(self, reference):
%         for i, ligne in enumerate(self._lignes):
%             if ligne[0] == reference:
%                 return i
%         return -1
%     def ajouter(self, reference, nombre):
%         assert nombre > 0
%         i = self._index(reference)
%         if i == -1:
%             self._lignes.append([reference, nombre])
%         else:
%             self._lignes[i][1] += nombre
%     def retirer(self, reference, nombre):
%         assert nombre > 0
%         i = self._index(reference)
%         assert i != -1 and nombre <= self._lignes[i][1]
%         self._lignes[i][1] -= nombre
%         if self._lignes[i][1] == 0:
%             del self._lignes[i]
%     def quantite(self, reference):
%         i = self._index(reference)
%         return 0 if i == -1 else self._lignes[i][1]
%     def references(self):
%         return sorted(ligne[0] for ligne in self._lignes)
%     def est_vide(self):
%         return self._lignes == []
% # Q4 : le meme scenario donne les memes reponses sur les deux implementations
% def scenario(fabrique):
%     stock = fabrique()
%     trace = [stock.est_vide()]
%     stock.ajouter("cle", 4); stock.ajouter("vis", 10); stock.ajouter("cle", 3)
%     trace.append((stock.quantite("cle"), stock.references()))
%     stock.retirer("vis", 10)
%     trace.append((stock.references(), stock.quantite("vis"), stock.est_vide()))
%     return trace
% assert scenario(Stock) == scenario(StockListe)
% # Q5 : le cout de `quantite` n'est pas le meme
% lectures = {"n": 0}
% class StockCompte(StockListe):
%     def _index(self, reference):
%         for i, ligne in enumerate(self._lignes):
%             lectures["n"] += 1
%             if ligne[0] == reference:
%                 return i
%         return -1
% s2 = StockCompte()
% for k in range(100):
%     s2.ajouter("ref%03d" % k, 1)
% lectures["n"] = 0
% s2.quantite("ref099")
% assert lectures["n"] == 100
% lectures["n"] = 0
% s2.quantite("ref000")
% assert lectures["n"] == 1
% END-VERIFY
\section*{Sujet S5 --- Exercice 1 : une classe pour un stock, interface et implémentation}

\textit{Durée conseillée : 70 minutes. Cet exercice est indépendant des deux
autres du sujet.}

\medskip
\textit{Contenu original Nexus --- ce n'est pas un sujet officiel d'examen.}

\subsection*{Énoncé}

Un atelier gère un stock de pièces. On spécifie une structure \lstinline{Stock}
par son \textbf{interface} seule :

\begin{center}
\begin{tabular}{|l|l|}
\hline
\textbf{Opération} & \textbf{Effet} \\
\hline
\texttt{ajouter(ref, n)} & ajoute $n$ pièces de la référence \\
\hline
\texttt{retirer(ref, n)} & retire $n$ pièces de la référence \\
\hline
\texttt{quantite(ref)} & quantité disponible, $0$ si absente \\
\hline
\texttt{references()} & liste triée des références en stock \\
\hline
\texttt{est\_vide()} & vrai si aucune pièce n'est en stock \\
\hline
\end{tabular}
\end{center}

Une référence dont la quantité tombe à zéro \textbf{disparaît} du stock : elle
n'apparaît plus dans \lstinline{references()}.

\begin{enumerate}
  \item Écrire une classe \lstinline{Stock} implémentée par un dictionnaire.
    Dérouler l'état après : ajouter $4$ clés, ajouter $10$ vis, ajouter $3$
    clés, retirer $10$ vis. Donner à chaque étape \lstinline{references()}.

  \item Écrire les préconditions de \lstinline{ajouter} et de
    \lstinline{retirer}, et les rendre exécutables. Donner cinq appels qui
    doivent être refusés, un par précondition violée. Vérifier qu'un appel
    refusé ne modifie pas le stock.

  \item Écrire une \textbf{seconde} implémentation \lstinline{StockListe} qui
    n'utilise aucun dictionnaire, mais une liste de couples
    $[\text{référence}, \text{nombre}]$. L'interface doit rester identique, au
    caractère près.

  \item Écrire une fonction \lstinline{scenario(fabrique)} qui joue la séquence
    de la question 1 sur une classe passée en paramètre et renvoie la trace des
    réponses. Vérifier qu'elle donne le même résultat pour les deux classes.
    Que prouve exactement cette égalité, et que ne prouve-t-elle pas ?

  \item On compte les lignes examinées par \lstinline{quantite} dans
    \lstinline{StockListe}, sur un stock de cent références. Donner ce nombre
    pour la dernière référence ajoutée, puis pour la première. Comparer avec
    l'implémentation par dictionnaire. Conclure sur ce que l'interface ne dit
    pas.
\end{enumerate}

\subsection*{Corrigé}

\textbf{Question 1.} \emph{(4 points)}

\begin{python}
class Stock:
    def __init__(self):
        self._quantites = {}
    def ajouter(self, reference, nombre):
        assert nombre > 0, "ajout non strictement positif"
        self._quantites[reference] = self._quantites.get(reference, 0) + nombre
    def retirer(self, reference, nombre):
        assert nombre > 0, "retrait non strictement positif"
        disponible = self._quantites.get(reference, 0)
        assert nombre <= disponible, "stock insuffisant"
        if nombre == disponible:
            del self._quantites[reference]
        else:
            self._quantites[reference] = disponible - nombre
    def quantite(self, reference):
        return self._quantites.get(reference, 0)
    def references(self):
        return sorted(self._quantites)
    def est_vide(self):
        return self._quantites == {}
\end{python}

\begin{center}
\begin{tabular}{|l|l|}
\hline
\textbf{Après} & \texttt{references()} \\
\hline
départ & $[\,]$ \\
\hline
$+4$ clés & $[\text{cle}]$ \\
\hline
$+10$ vis & $[\text{cle}, \text{vis}]$ \\
\hline
$+3$ clés & $[\text{cle}, \text{vis}]$, quantité clé $= 7$ \\
\hline
$-10$ vis & $[\text{cle}]$ \\
\hline
\end{tabular}
\end{center}

\medskip
\textbf{Question 2.} \emph{(5 points)}

Préconditions : le nombre ajouté ou retiré est \textbf{strictement positif} ;
un retrait ne peut excéder la quantité disponible, et une référence absente a
une quantité de $0$.

Cinq appels refusés : \lstinline{ajouter("cle", 0)},
\lstinline{ajouter("cle", -1)}, \lstinline{retirer("cle", 0)},
\lstinline{retirer("cle", 99)}, \lstinline{retirer("absent", 1)}.

Chaque assertion est placée \textbf{avant} toute modification : l'état du stock
est donc inchangé après un refus.

\medskip
\textbf{Question 3.} \emph{(5 points)}

\begin{python}
class StockListe:
    def __init__(self):
        self._lignes = []
    def _index(self, reference):
        for i, ligne in enumerate(self._lignes):
            if ligne[0] == reference:
                return i
        return -1
    def ajouter(self, reference, nombre):
        assert nombre > 0
        i = self._index(reference)
        if i == -1:
            self._lignes.append([reference, nombre])
        else:
            self._lignes[i][1] += nombre
    def retirer(self, reference, nombre):
        assert nombre > 0
        i = self._index(reference)
        assert i != -1 and nombre <= self._lignes[i][1]
        self._lignes[i][1] -= nombre
        if self._lignes[i][1] == 0:
            del self._lignes[i]
    def quantite(self, reference):
        i = self._index(reference)
        return 0 if i == -1 else self._lignes[i][1]
    def references(self):
        return sorted(ligne[0] for ligne in self._lignes)
    def est_vide(self):
        return self._lignes == []
\end{python}

\medskip
\textbf{Question 4.} \emph{(4 points)}

\lstinline{scenario(Stock) == scenario(StockListe)}.

L'égalité prouve que les deux implémentations sont \textbf{indiscernables sur
ce scénario} : un utilisateur qui ne passe que par l'interface ne peut pas
dire laquelle il manipule. Elle ne prouve \textbf{pas} qu'elles se comportent
identiquement sur tous les scénarios, ni qu'elles ont le même coût.

\medskip
\textbf{Question 5.} \emph{(2 points)}

Dernière référence ajoutée : $\mathbf{100}$ lignes examinées. Première :
$\mathbf{1}$. Avec le dictionnaire, l'accès ne dépend pas de la position ni du
nombre de références.

L'interface décrit ce que les opérations \textbf{font}, jamais ce qu'elles
\textbf{coûtent}. Deux implémentations conformes peuvent être également
correctes et très inégalement utilisables.

\subsection*{Critique}

\begin{itemize}
  \item \textbf{Oublier la disparition à zéro.} C'est une exigence de l'énoncé,
    pas un détail : une implémentation qui laisse une ligne à $0$ échoue au
    scénario de la question 4 sur \lstinline{references()}, et l'échec ne dit
    pas laquelle des deux classes est fautive.
  \item \textbf{Assertion après modification.} Placer le contrôle de stock après
    la soustraction laisse une quantité négative en mémoire avant d'échouer. La
    question 2 exige explicitement de vérifier l'état après un refus.
  \item \textbf{« Le test passe, donc les deux classes sont équivalentes ».} La
    question 4 demande la limite de cette conclusion, et c'est là que se joue la
    différence entre un candidat qui teste et un candidat qui sait ce qu'un test
    établit.
  \item \textbf{Accéder à \lstinline{_lignes} depuis \lstinline{scenario}.} Le
    scénario ne doit utiliser que l'interface, sinon il cesse de valoir pour les
    deux classes.
\end{itemize}

\subsection*{Portée pédagogique}

L'exercice fait écrire \textbf{deux} fois la même structure pour rendre visible
ce qui, autrement, reste une phrase de cours : l'interface est ce qui survit au
changement d'implémentation.

La question 4 en donne la vérification pratique — un même scénario, deux
classes, une même trace — et la question 5 en donne la limite, sans laquelle
l'idée serait fausse dans l'autre sens : rien, dans une interface, n'oblige
deux implémentations à être également efficaces.
